% Section 1 of the manuscript: introduction. It follows the section plan and the contributions of outline.md; the
% claims it makes are those of the claims table, stated with the sections that give their evidence.

\section{Introduction}
\label{sec:introduction}

Keystroke logging records each key press of a writer with its time, so that the process of writing can be studied
alongside the text it produces \citep{wengelin2023}. From such logs researchers compute inter-key intervals
and pauses, revisions, bursts of production between pauses, and the amount of text produced and the time taken
\citep{conijn2019}.

When Japanese is typed on a keyboard, the text passes through an input method: the writer types a reading, in Roman
letters or in kana, the input method converts it into kanji and kana, and the writer accepts the conversion, changes it
or abandons it. The conversion is part of the process of writing. Loggers of writing handle it in three ways
(Section~\ref{sec:related}). A logger that keeps only the final text does not see it. A logger of keys does not record
it directly: the conversion can be inferred from the presses by rules whose validity has to be shown, and the log holds
the keys that spell the reading. A logger that observes the conversion reads the composition string or the text on the
screen.

A logger that records the conversion, directly or by inference, therefore also records the text that a writer produces
or the keys that spell it. A log that holds neither the text nor the names of keys changes what has to be protected
when the log is kept or shared, but only if it is known what the log records and how much of the text it still reveals.
This article asks three questions: what can be logged about the conversion without reading text; how closely such a log
agrees with the sessions that an input field receives from the input method; and how much of the text the log still
reveals.

The article makes five contributions.
\begin{enumerate}
\item An analysis of observability (Section~\ref{sec:observability}): what the documented public interfaces of
Windows and macOS let a program that reads no text learn about the conversion in another application. Such a program
may receive key codes, which spell the reading when it is typed in Roman letters, but uses no interface that delivers
the composition, the candidates or other text; the logger writes only the category of each key. On both systems an
application's own accessibility implementation can report a conversion without text, but whether it does depends on the
application. Among the routes that do not depend on the application, Windows reports the display and the cloaking of
the input method's windows without text, and on macOS none of the documented interfaces examined reports a conversion,
so the logger infers it from key presses.
\item A logger for both systems and its validation (Sections~\ref{sec:logger} to \ref{sec:results-text}). The logger
writes a category and a time for every press and release of the keys it observes and one summary row for every
conversion session, and it declares in its log which rows the log can contain. Its rows were compared with the
sessions assembled from the calls that a reference input field received from the input method, in three cells (macOS
with live conversion on and off; Windows with the Microsoft input method), under a pre-registered protocol with
synthetic input. On Windows, where the logger records the windows of the input method and the reference the
composition, the agreement includes the difference between the two (Section~\ref{sec:observability-consequence}). The
estimate of this agreement is the primary contribution of the article.
\item An estimate of text recoverability (Section~\ref{sec:recoverability}): how accurately an attacker can identify
the typed phrase among $K$ candidates, or reconstruct the typed reading, from the key rows of a log, and whether either
of two granularities of the key rows that keep no time keeps the identification near chance under the registered rule.
\item The runtime cost of the logger on each system, with the machine idle and with scripted typing
(Section~\ref{sec:cost}).
\item A dataset (Section~\ref{sec:dataset}): the logs of the measurement with the ground truth of the reference input
field, the plans and the record of every press sent, the phrase corpus and the runtime samples. With these records
another rule that infers sessions from key presses can be evaluated against the same ground truth without new
measurement.
\end{enumerate}

The analyses of agreement and of recoverability come from the same measurement runs; the runtime cost was measured in
runs of its own, and the dataset holds the runs of the calibration pilots as a separate part. All key input of the
measurement runs was produced by a program, and no claim is made about natural typing, about other input methods or
about applications other than the reference input fields (Section~\ref{sec:limitations}).
